
Benchmark contamination evades standard detection when training data is paraphrased. Models trained on rephrased MMLU items achieve 85.9\% accuracy while bypassing n-gram filters \citep{yang2023rethinking}. This work investigates whether contamination leaves behavioral signatures detectable via confidence patterns. The Semantic Saturation Index (SSI) is proposed, measuring inverse confidence variance across paraphrases of benchmark items, hypothesizing that contaminated models exhibit uniform confidence due to robust semantic representations. Through controlled experiments on Mistral-7B (0--50\% MMLU contamination, K=20 paraphrases), the causal mechanism is validated: contamination creates training exposure (31.1\% accuracy gain at 50\% contamination), paraphrase-augmented training produces representation invariance (Mean Pairwise Similarity difference 0.065, Cohen's d=0.52, p=0.008), and invariance correlates with confidence uniformity (Pearson r=$-0.517)$. However, the SSI metric fails on real inference data, achieving AUC=0.506 (chance level). The inverse-variance formulation amplifies noise: within-group standard deviation exceeds mean SSI at all contamination levels. This work provides the first empirical demonstration of the contamination-to-uniformity mechanism chain, documents the failure of inverse-variance metrics, and identifies a simulation-reality gap where simulated experiments passed but real-data evaluation revealed metric failure.
